\chapter{CASOS DE PRUEBA FUNCIONALES}
\label{anx:casos}

Los casos de prueba se ejecutan sobre el \emph{development build} instalado
en un dispositivo Android, con una cuenta de prueba y el catálogo de máquinas
cargado en Firestore. En la columna \emph{Resultado} se registra si el caso
pasa o falla.

% TODO (autora): al ejecutar cada caso, reemplazar \cpres por "Pasa" o
% "Falla (ver observación)". Para marcar todos como pendientes en gris:
\newcommand{\cpres}{\textcolor{black!45}{\footnotesize\shortstack[l]{$\square$ Pasa\\$\square$ Falla}}}

{\small
\begin{longtable}{|L{1.15cm}|L{1.0cm}|L{5.2cm}|L{4.4cm}|C{1.5cm}|}
\caption{Casos de prueba funcionales}
\label{tab:casos-prueba}\\
\hline
\rowcolor{encabezado} \textbf{ID} & \textbf{RF} & \textbf{Pasos} & \textbf{Resultado esperado} &
\textbf{Resultado} \\
\hline
\endfirsthead
\hline
\rowcolor{encabezado} \textbf{ID} & \textbf{RF} & \textbf{Pasos} & \textbf{Resultado esperado} &
\textbf{Resultado} \\
\hline
\endhead

\endfoot
\multicolumn{5}{|l|}{\textbf{Autenticación, onboarding y perfil}}\\ \hline
CP-01 & RF1 & Registrarse con nombre, apellido, correo válido y contraseña
de 6 o más caracteres. & Se crea la cuenta y el perfil; se abre el
onboarding. & \cpres \\ \hline
CP-02 & RF1 & Registrarse con contraseñas distintas o de menos de 6
caracteres. & Se muestra el mensaje de validación; no se crea la cuenta. &
\cpres \\ \hline
CP-03 & RF1 & Iniciar sesión con Google por primera vez y luego cerrar
sesión y volver a ingresar. & El perfil se crea solo la primera vez; no se
sobrescribe. & \cpres \\ \hline
CP-04 & RF2 & Ingresar con un usuario sin objetivo ni nivel. & Se muestra el
onboarding; las pestañas de la app no están disponibles. & \cpres \\ \hline
CP-05 & RF2 & Elegir objetivo, nivel y ``Generar con IA''. & Se guardan
objetivo y nivel; se abre la generación de rutina. & \cpres \\ \hline
CP-06 & RF1 & Solicitar la recuperación de contraseña. & Llega el correo de
restablecimiento. & \cpres \\ \hline
CP-07 & RF12 & Editar la estatura y el objetivo en el perfil; cerrar
sesión. & Los cambios persisten; al cerrar sesión se vuelve a la
bienvenida. & \cpres \\
\hline
\multicolumn{5}{|l|}{\textbf{Guía de máquinas y ejercicios}}\\ \hline
CP-08 & RF3 & Abrir la pestaña Máquinas. & Se muestran las 35 máquinas en
cuadrícula. & \cpres \\ \hline
CP-09 & RF3 & Buscar ``biceps'' y luego ``bíceps''. & Ambas búsquedas
devuelven los mismos resultados. & \cpres \\ \hline
CP-10 & RF3 & Aplicar un filtro de grupo muscular y uno de tipo. & Solo se
muestran las máquinas que cumplen ambos filtros. & \cpres \\ \hline
CP-11 & RF4 & Abrir la ficha de la prensa de piernas. & Se ven descripción,
músculos, guía, errores, seguridad y variantes. & \cpres \\ \hline
CP-12 & RF5 & En la ficha, revisar la sección de ejercicios. & Cada
ejercicio muestra su animación, músculo objetivo y dificultad. & \cpres \\ \hline
CP-13 & RF5 & Abrir una ficha sin conexión a Internet tras haberla visto
antes. & Se muestra el contenido en caché; no se bloquea la pantalla. &
\cpres \\
\hline
\multicolumn{5}{|l|}{\textbf{Escáner con realidad aumentada}}\\ \hline
CP-14 & RF6 & Abrir el escáner por primera vez. & Se solicita el permiso de
cámara con el mensaje explicativo. & \cpres \\ \hline
CP-15 & RF6 & Enfocar una máquina y capturar. & Al no haber modelo, se
muestra la lista para confirmar la máquina. & \cpres \\ \hline
CP-16 & RF6 & Elegir una máquina de la lista. & Su ficha resumida se
superpone sobre la cámara en vivo. & \cpres \\ \hline
CP-17 & RF6 & Desde una ficha, tocar ``Ver en Realidad Aumentada''. & El
escáner abre con esa máquina ya superpuesta. & \cpres \\
\hline
\multicolumn{5}{|l|}{\textbf{Rutinas y planificación}}\\ \hline
CP-18 & RF7 & Generar una rutina con objetivo ``Ganar masa'' y nivel
``Principiante''. & Rutina de 3 días (empuje, tirón, pierna) con 4 series de
9 repeticiones y 90 s de descanso. & \cpres \\ \hline
CP-19 & RF7 & Generar una rutina con ``Perder peso'' e ``Intermedio''. &
Rutina de 4 días con un día de cardio y 3 series de 15 repeticiones. &
\cpres \\ \hline
CP-20 & RF8 & Crear una rutina personalizada de 2 días con ejercicios. & La
rutina se guarda y queda como activa. & \cpres \\ \hline
CP-21 & RF8 & Intentar guardar un día sin ejercicios. & Se muestra el motivo
y no se guarda. & \cpres \\ \hline
CP-22 & RF8 & Desde una ficha, agregar un ejercicio al día 1; repetir. & La
primera vez se agrega; la segunda se informa que ya existe. & \cpres \\ \hline
CP-23 & RF8 & Activar otra rutina y eliminar la anterior. & El inicio muestra
la nueva rutina activa; la eliminada desaparece. & \cpres \\
\hline
\multicolumn{5}{|l|}{\textbf{Nutrición}}\\ \hline
CP-24 & RF9 & Calcular el IMC con 70~kg y 175~cm. & IMC 22.9, categoría
normal. & \cpres \\ \hline
CP-25 & RF9 & Calcular con una estatura de 30~cm. & Mensaje de validación;
no se guarda. & \cpres \\ \hline
CP-26 & RF9 & Revisar Firestore tras un cálculo. & Existe el registro con la
fecha del servidor. & \cpres \\
\hline
\multicolumn{5}{|l|}{\textbf{Progreso y reportes}}\\ \hline
CP-27 & RF10 & Abrir Progreso con estatura en el perfil. & La estatura
aparece precargada. & \cpres \\ \hline
CP-28 & RF10 & Registrar 80~kg y luego 77.5~kg. & Se crean dos registros con
su IMC. & \cpres \\ \hline
CP-29 & RF11 & Observar el reporte tras el caso anterior. & El gráfico se
actualiza al instante y la variación es $-2.5$~kg. & \cpres \\ \hline
CP-30 & RF11 & Registrar diez pesos distintos. & El gráfico muestra los ocho
últimos. & \cpres \\ \hline
\end{longtable}
}

\section{Script de prueba de la lógica de negocio}

Las pruebas de la sección~\ref{sec:pruebas} sobre el generador de rutinas y
el cálculo del IMC se ejecutaron con el siguiente \emph{script}, usando
Node.js~22 con soporte nativo de tipos (\archivo{node
--experimental-strip-types prueba.ts}).

\begin{lstlisting}[language=TypeScript, caption={Script de prueba de las funciones puras}, label={lst:script-prueba}]
import { generarPlanHeuristico, repartirCantidad } from "./generadorRutinas.ts";
import { calcularIMC, clasificarIMC } from "./imcCalculator.ts";

const objetivos = ["ganar_masa", "mantenimiento", "perder_peso", "resistencia"] as const;
const niveles = ["principiante", "intermedio", "avanzado"] as const;

for (const o of objetivos) for (const n of niveles) {
  const plan = generarPlanHeuristico(o, n);
  console.log(o, n, plan.diasPorSemana, plan.parametrosSerie,
    plan.bloques.map((b) => b.enfoque).join(" | "),
    plan.bloques.reduce((total, b) => total + b.cantidadEjercicios, 0));
}
console.log(repartirCantidad(5, 3), repartirCantidad(6, 3), repartirCantidad(4, 2));

for (const [peso, altura] of [[50, 170], [53.4, 170], [70, 175],
                              [72.2, 170], [86.7, 170], [100, 170]]) {
  const imc = calcularIMC(peso, altura);
  console.log(peso, altura, imc, clasificarIMC(imc));
}
for (const v of [18.4, 18.5, 24.9, 25, 29.9, 30]) console.log(v, clasificarIMC(v));
\end{lstlisting}
